\section{Discussão e Conclusão}

Os resultados obtidos ao longo deste trabalho experimental revelam uma elevada concordância entre os valores experimentais e teóricos, quer na determinação da resistência de entrada de um circuito como na resistência de saída de uma fonte de tensão real.

Esta experiência permitiu assim validar o Teorema de Thévenin, uma vez que os resultados obtidos foram exatos (com exatidões de 99,36\% para o caso a) e 99,48\% para o caso b) e 99,07\% para a montagem 2).

Possiveís erros experimentais que podem ter afastado os valores experimetais dos valores teóricos são erros de leitura dos instrumentos, imperfeições das ligações do circuito e tolerâncias dos valores associados às resistências. Além disso, as pequenas discrepâncias são justificadas pela incerteza de leitura dos instrumentos de medição, uma vez que o voltímetro e o amperímetro possuem uma resistência interna que podem introduzir imprecisões nas medições de corrente e tensão; pelo próprio aquecimento resistivo (efeito Joule) durante a passagem da corrente.